\documentclass[12pt]{article}
\usepackage[utf8]{inputenc}
\usepackage[T1]{fontenc}
\usepackage{hyperref}
\title{Binukid Language Resources for K-12 Curriculum Development}
\author{Compiled by Claude Sonnet 5.5 (claude-sonnet-5-5) for Andr\'as Kornai}
\date{}
\begin{document}
\maketitle

\section{Introduction}
Binukid (Bukid, ISO 639-3 \texttt{bkd}) is an Austronesian language of Bukidnon province, Northern Mindanao, Philippines. The Malaybalay variety is regarded as the prestige variety. Documentation found is thin; phonology is the best-covered area.

\section{Primary Educational Resources}
No primers, textbooks or school programmes were found.

\section{Linguistic and Reference Materials}
\begin{itemize}
\item Wikipedia overview: \url{https://en.wikipedia.org/wiki/Bukid_language}
\item Hermocilla-Borres (2018), ``Phonological Features of Contemporary Spoken Binukid'', Journal of Modern Languages: \url{https://jml.um.edu.my/article/view/15611}
\end{itemize}
No full grammar was found.

\section{Religious and Community Resources}
None found.

\section{Digital and Audio Resources}
\begin{itemize}
\item Glosbe community-built dictionary (user-contributed, reliability varies): \url{https://glosbe.com/bkd/lib}
\item Wikitongues language profile: \url{https://wikitongues.org/languages/bkd/}
\end{itemize}

\section{Curriculum Development Context}
A curriculum would need substantial new lexical and grammatical documentation. The 2018 phonology article and the SIL phoneme paper by Atherton (1953, mentioned in search results but without a URL seen) are starting points for an orthography.
\end{document}